\documentclass[12pt,a4paper]{article}
\usepackage[utf8]{inputenc}
\usepackage[spanish]{babel}
\usepackage{graphicx}
\usepackage{xcolor}
\usepackage{listings}
\usepackage{hyperref}
\usepackage{geometry}
\usepackage{tikz}
\usetikzlibrary{positioning, arrows.meta, shapes.geometric}

\geometry{margin=2.5cm}

% Configuración de colores
\definecolor{firebaseYellow}{RGB}{255,202,40}
\definecolor{firebaseOrange}{RGB}{255,107,0}
\definecolor{codebg}{RGB}{245,245,245}

% Configuración de código
\lstset{
    backgroundcolor=\color{codebg},
    basicstyle=\ttfamily\small,
    breaklines=true,
    frame=single,
    numbers=left,
    numberstyle=\tiny,
    language=JavaScript
}

\title{\textbf{Guía de Arquitectura para Aplicación con Firebase}\\
\large Proyecto: PracticaPIL2G4}
\author{Diego - Guía Técnica}
\date{\today}

\begin{document}

\maketitle
\tableofcontents
\newpage

\section{Introducción}

Esta guía documenta la arquitectura de una aplicación web que utiliza Firebase como backend. El proyecto está desarrollado en Ubuntu y sigue el patrón MVC (Modelo-Vista-Controlador) adaptado para trabajar con servicios en la nube.

\subsection{Objetivos de la Arquitectura}
\begin{itemize}
    \item Separación clara de responsabilidades (MVC)
    \item Integración nativa con servicios Firebase
    \item Flujo de desarrollo local con despliegue automatizado
    \item Escalabilidad y mantenibilidad del código
\end{itemize}

\section{Componentes de Firebase}

\subsection{Firestore (Base de Datos)}
Firestore es una base de datos NoSQL documental en tiempo real. En nuestra arquitectura:
\begin{itemize}
    \item \textbf{Colecciones}: Agrupaciones de documentos (ej: \texttt{eventos}, \texttt{usuarios})
    \item \textbf{Documentos}: Unidades de datos en formato JSON
    \item \textbf{Reglas de seguridad}: Definidas en \texttt{firestore.rules}
    \item \textbf{Índices}: Optimización de consultas en \texttt{firestore.indexes.json}
\end{itemize}

\subsection{Hosting}
Servicio que aloja los archivos estáticos de la aplicación:
\begin{itemize}
    \item Directorio raíz: \texttt{/public}
    \item Despliegue automático con \texttt{firebase deploy}
    \item SSL gratuito y CDN global
\end{itemize}

\subsection{Authentication}
Sistema de autenticación de usuarios:
\begin{itemize}
    \item Proveedores: Email/contraseña, Google, GitHub, etc.
    \item Integración con Firestore para perfiles de usuario
    \item Tokens JWT para sesiones seguras
\end{itemize}

\section{Arquitectura del Proyecto en Ubuntu}

\subsection{Estructura de Directorios}
\begin{lstlisting}[frame=none, numbers=none]
PracticaPIL2G4/
├── public/                 # Raíz web (Hosting)
│   ├── index.html          # Página principal
│   ├── css/                 # Estilos
│   │   └── style.css
│   ├── js/                  # JavaScript
│   │   ├── app.js           # Punto de entrada
│   │   ├── model/           # Capa Modelo
│   │   │   ├── Evento.js
│   │   │   └── Usuario.js
│   │   ├── view/            # Capa Vista
│   │   │   ├── renderizador.js
│   │   │   └── componentes.js
│   │   └── controller/       # Capa Controlador
│   │       ├── eventoController.js
│   │       └── authController.js
│   └── assets/              # Imágenes, fuentes, etc.
├── firebase.json            # Configuración Firebase
├── .firebaserc              # Alias de proyectos
├── firestore.rules          # Reglas de seguridad
├── firestore.indexes.json   # Índices de base de datos
├── .github/                 # GitHub Actions
│   └── workflows/
│       ├── firebase-hosting-merge.yml
│       └── firebase-hosting-pull-request.yml
├── model/                   # (Histórico - migrar a public/js/model)
├── view/                    # (Histórico - migrar a public/js/view)
├── javascript/              # (Histórico - migrar a public/js)
├── css/                     # (Histórico - migrar a public/css)
└── html/                    # (Histórico - migrar a public/)
\end{lstlisting}

\subsection{Flujo de Datos en Ubuntu}
\begin{tikzpicture}[node distance=2cm, auto, >=Stealth]
    % Nodos
    \node[rectangle, draw, fill=blue!20, rounded corners] (user) {Usuario};
    \node[rectangle, draw, fill=green!20, rounded corners, below=of user] (browser) {Navegador};
    \node[rectangle, draw, fill=yellow!20, rounded corners, below=of browser] (local) {Servidor Local (firebase serve)};
    \node[rectangle, draw, fill=orange!20, rounded corners, right=of local, xshift=3cm] (firebase) {Firebase Cloud};
    \node[rectangle, draw, fill=red!20, rounded corners, below=of firebase] (github) {GitHub};
    
    % Flechas
    \draw[->] (user) -- (browser);
    \draw[->] (browser) -- (local) node[midway, left] {localhost:5000};
    \draw[<->] (local) -- (firebase) node[midway, above] {Emuladores};
    \draw[->] (local) -- (github) node[midway, below] {git push};
    \draw[->] (github) -- (firebase) node[midway, right] {GitHub Action};
\end{tikzpicture}

\section{Capa Modelo: Conexión con Firebase}

\subsection{Configuración Inicial de Firebase}
Archivo: \texttt{public/js/conector/firebaseConfig.js}

\begin{lstlisting}
// Importar SDKs necesarios
import { initializeApp } from "firebase/app";
import { getFirestore } from "firebase/firestore";
import { getAuth } from "firebase/auth";

// Configuración del proyecto (obtenida de Firebase Console)
const firebaseConfig = {
    apiKey: "tu-api-key",
    authDomain: "pi-l2-g4.firebaseapp.com",
    projectId: "pi-l2-g4",
    storageBucket: "pi-l2-g4.appspot.com",
    messagingSenderId: "xxx",
    appId: "xxx"
};

// Inicializar Firebase
const app = initializeApp(firebaseConfig);
export const db = getFirestore(app);
export const auth = getAuth(app);
\end{lstlisting}

\subsection{Modelo de Eventos}
Archivo: \texttt{public/js/model/Evento.js}

\begin{lstlisting}
import { db } from '../conector/firebaseConfig.js';
import { collection, addDoc, getDocs, updateDoc, deleteDoc, doc } from "firebase/firestore";

export class Evento {
    constructor(titulo, fecha, descripcion) {
        this.titulo = titulo;
        this.fecha = fecha;
        this.descripcion = descripcion;
    }
    
    // Guardar en Firestore
    async guardar() {
        try {
            const docRef = await addDoc(collection(db, "eventos"), {
                titulo: this.titulo,
                fecha: this.fecha,
                descripcion: this.descripcion,
                timestamp: new Date()
            });
            return docRef.id;
        } catch (error) {
            console.error("Error guardando evento:", error);
            throw error;
        }
    }
    
    // Obtener todos los eventos
    static async obtenerTodos() {
        const querySnapshot = await getDocs(collection(db, "eventos"));
        return querySnapshot.docs.map(doc => ({
            id: doc.id,
            ...doc.data()
        }));
    }
}
\end{lstlisting}

\section{Capa Controlador: Lógica de Negocio}

\subsection{Controlador de Eventos}
Archivo: \texttt{public/js/controller/eventoController.js}

\begin{lstlisting}
import { Evento } from '../model/Evento.js';
import { renderizarEventos } from '../view/renderizador.js';

export class EventoController {
    constructor() {
        this.eventos = [];
    }
    
    async cargarEventos() {
        try {
            this.eventos = await Evento.obtenerTodos();
            renderizarEventos(this.eventos);
        } catch (error) {
            console.error("Error cargando eventos:", error);
            this.mostrarError("No se pudieron cargar los eventos");
        }
    }
    
    async crearEvento(datosEvento) {
        const evento = new Evento(
            datosEvento.titulo,
            datosEvento.fecha,
            datosEvento.descripcion
        );
        
        try {
            await evento.guardar();
            await this.cargarEventos(); // Recargar lista
        } catch (error) {
            console.error("Error creando evento:", error);
        }
    }
}
\end{lstlisting}

\section{Capa Vista: Renderizado en el Navegador}

\subsection{Renderizador de Eventos}
Archivo: \texttt{public/js/view/renderizador.js}

\begin{lstlisting}
export function renderizarEventos(eventos) {
    const contenedor = document.getElementById('lista-eventos');
    
    if (!contenedor) return;
    
    contenedor.innerHTML = eventos.map(evento => `
        <div class="evento-card" data-id="${evento.id}">
            <h3>${evento.titulo}</h3>
            <p>Fecha: ${new Date(evento.fecha).toLocaleDateString()}</p>
            <p>${evento.descripcion}</p>
            <button onclick="eliminarEvento('${evento.id}')">Eliminar</button>
        </div>
    `).join('');
}
\end{lstlisting}

\section{Autenticación de Usuarios}

\subsection{Controlador de Autenticación}
Archivo: \texttt{public/js/controller/authController.js}

\begin{lstlisting}
import { auth } from '../conector/firebaseConfig.js';
import { 
    signInWithEmailAndPassword,
    createUserWithEmailAndPassword,
    signOut,
    onAuthStateChanged 
} from "firebase/auth";

export class AuthController {
    constructor() {
        this.usuarioActual = null;
        this.observarEstadoAuth();
    }
    
    async login(email, password) {
        try {
            const userCredential = await signInWithEmailAndPassword(auth, email, password);
            this.usuarioActual = userCredential.user;
            return true;
        } catch (error) {
            console.error("Error login:", error);
            return false;
        }
    }
    
    async registrar(email, password) {
        try {
            const userCredential = await createUserWithEmailAndPassword(auth, email, password);
            return userCredential.user;
        } catch (error) {
            console.error("Error registro:", error);
            throw error;
        }
    }
    
    observarEstadoAuth() {
        onAuthStateChanged(auth, (user) => {
            this.usuarioActual = user;
            this.actualizarUI(user);
        });
    }
    
    actualizarUI(user) {
        const loginSection = document.getElementById('login-section');
        const appSection = document.getElementById('app-section');
        
        if (user) {
            loginSection.style.display = 'none';
            appSection.style.display = 'block';
        } else {
            loginSection.style.display = 'block';
            appSection.style.display = 'none';
        }
    }
}
\end{lstlisting}

\section{Reglas de Seguridad en Firestore}

\subsection{firestore.rules}
\begin{lstlisting}
rules_version = '2';
service cloud.firestore {
    match /databases/{database}/documents {
        // Reglas para eventos
        match /eventos/{eventoId} {
            allow read: if true;  // Cualquiera puede leer
            allow write: if request.auth != null;  // Solo usuarios autenticados
        }
        
        // Reglas para usuarios
        match /usuarios/{userId} {
            allow read: if true;
            allow write: if request.auth != null && request.auth.uid == userId;
        }
    }
}
\end{lstlisting}

\section{Despliegue Automatizado con GitHub Actions}

\subsection{Workflow de Despliegue}
Archivo: \texttt{.github/workflows/firebase-hosting-merge.yml}

\begin{lstlisting}[language=yaml]
name: Deploy to Firebase Hosting on merge
'on':
  push:
    branches:
      - main
jobs:
  build_and_deploy:
    runs-on: ubuntu-latest
    steps:
      - uses: actions/checkout@v3
      - run: npm ci && npm run build
      - uses: FirebaseExtended/action-hosting-deploy@v0
        with:
          repoToken: '${{ secrets.GITHUB_TOKEN }}'
          firebaseServiceAccount: '${{ secrets.FIREBASE_SERVICE_ACCOUNT_PI_L2_G4 }}'
          channelId: live
          projectId: pi-l2-g4
\end{lstlisting}

\section{Comandos Útiles en Ubuntu}

\subsection{Desarrollo Local}
\begin{lstlisting}[frame=none, numbers=none]
# Iniciar emuladores de Firebase
firebase emulators:start

# Servir solo hosting local
firebase serve

# Servir en puerto específico
firebase serve --port 3000
\end{lstlisting}

\subsection{Despliegue}
\begin{lstlisting}[frame=none, numbers=none]
# Desplegar todo (hosting + reglas + índices)
firebase deploy

# Desplegar solo hosting
firebase deploy --only hosting

# Desplegar solo reglas de Firestore
firebase deploy --only firestore:rules
\end{lstlisting}

\section{Conclusión}

Esta arquitectura proporciona:

\begin{itemize}
    \item \textbf{Separación clara}: MVC bien definido con Firebase como backend
    \item \textbf{Flujo de trabajo profesional}: Desarrollo local con emuladores, despliegue automatizado
    \item \textbf{Seguridad}: Reglas de Firestore y autenticación integrada
    \item \textbf{Escalabilidad}: Firebase maneja automáticamente la infraestructura
\end{itemize}

La migración de los archivos estáticos a la carpeta \texttt{public} es el único paso necesario para que esta arquitectura funcione correctamente en producción.

\end{document}